\section{Relation to numerical automata and positive realization}\label{sec:comparison54}
The input/output object and the order of quantifiers are decisive in comparing realization theorems. Classical probabilistic automata assign numerical acceptance probabilities to words, while stochastic-language results also study threshold recognition \cite{Rabin,Paz}. A numerical equivalence of two fixed machines is not the assertion that for every horizon there exists a new bounded-width approximate machine. Conversely, language equivalence need not preserve the numerical probabilities.

For weighted automata over a field, the finite-Hankel-rank characterization gives finite linear representations and identifies their minimum dimension \cite{BR}. It does not enforce nonnegative rows, stochastic normalization, or simplex-valued decoders. Our common-row contraction concerns these constraints simultaneously. The invariant observable quotient in Theorem~\ref{thm:observable53} uses the same reachable/invisible linear reduction; its additional conclusion is a finite-group criterion for exact positive realization with the full horizon-dependent quantifier.

The classical positive-realization problem asks for a positive system representing specified input/output data, and distinguishes positive dimension from unrestricted linear dimension \cite{BF04}. Invariant cones and polyhedral sections are part of that theory, not new notions introduced here. Proposition~\ref{prop:polygon54} is an explicit common invariant-enclosure construction. The new matched statement pairs it with a full-subcritical lower estimate against all clocked stochastic machines. Neither one cutwise nonnegative factorization nor a finite-rank Hankel matrix supplies that lower estimate.

Brodsky and Pippenger identify bounded-error measure-once quantum languages as group languages \cite[Theorem 3.3]{BP}. Their probabilistic simulation result preserves a cutpoint language with a potentially changed cutpoint \cite[Theorem 3.6]{BP}; it is not uniform additive preservation of every acceptance probability. Our categorical unitary corollary fixes an entire terminal measurement law in total variation and allows horizon-dependent stochastic approximants. The distinction is already visible in the $d$-outcome threshold $1-1/d$.

The closure theorem and strict-threshold decision method of Blondel, Jeandel, Koiran, and Portier \cite{BJKP}, and the effective Zariski-closure algorithms \cite{DJK,NPSW}, are algorithmic premises of Section~\ref{sec:effective54}. Computing the closure is not a new claim of this article. The additional outputs are a constrained minimax noise threshold, an attaining categorical decoder, and a terminating construction of an all-machine occupation constant. Recent Lie-algebra-based closure computations \cite{deGraaf} give an alternative representation of the algebraic closure, but their complex connected components still require real-component analysis for the present application.

Peter--Weyl approximation, Haar averaging, and the representative-function algebra are classical compact-group and almost-periodic-function machinery \cite{Folland}. Uniform approximation by a finite sum of matrix coefficients gives a finite \emph{linear} representation of an output function. It does not give positive stochastic rows or the repeated information loss quantified by \eqref{eq:occupation54}. Similarly, Steinberg's minimal compact topological automaton and enveloping-monoid results \cite[Theorems 2.2 and 2.6]{Steinberg} concern recognition of a language. His compact topological-monoid criterion \cite[Proposition 2.8]{Steinberg} uses clopen recognition, not a continuous simplex-valued response with a prescribed error norm. Our return condition does not classify arbitrary irreversible semigroups.

Distribution-based bisimulation metrics provide robust comparisons for specified probabilistic automata \cite{FSZ}. They do not by themselves impose a minimum over all horizon-specific stochastic realizations of an external group experiment. The finite future-output equivalence in Theorem~\ref{thm:quotient54} is deliberately a component-quotient statement, rather than a claim that approximate equivalence always has a unique minimal positive stochastic realization.

\begin{center}\small
\begin{tabular}{p{.27\textwidth}p{.31\textwidth}p{.32\textwidth}}
\toprule
Comparison object & Error or equivalence & Machine quantifier\\
\midrule
Weighted series \cite{BR} & Exact scalar coefficients; linear rank & One linear representation\\
Group-language simulation \cite{BP} & Cutpoint language; not additive probability error & One automaton\\
Positive realization \cite{BF04} & Exact specified numerical response & One positive system\\
Topological recognition \cite{Steinberg} & Clopen language acceptance & One recognizing action\\
Bisimulation metrics \cite{FSZ} & Distributional comparison of given systems & Two specified automata\\
Theorem~\ref{thm:radius54} & Worst-word norm error; TV for a whole categorical law & A new clocked stochastic machine at every horizon\\
Theorem~\ref{thm:matched54} & Binary TV, every fixed $\varepsilon<\rho/2$ & Same nonuniform lower model; exact upper construction\\
\bottomrule
\end{tabular}
\end{center}
The contribution is the joint conclusion of a sharp constrained radius, finite permutation attainment, arbitrary-width occupation below the threshold, and, for the specified Diophantine family, a matching full-subcritical growth law. The elementary radius center, the arithmetic pigeonhole argument, the closure algorithm, and the harmonic-analysis ingredients are not asserted individually new. The accompanying theorem-level audit records the closest comparisons and their changed hypotheses; it is not a claim of exhaustive priority clearance.

\subsection{Additional mathematical material}
The preceding article and its complete supplement remain separately buildable. They retain the scalar two-extreme proof, compatible reachable sections, exact stationarization, one-surplus geometry, stable sections, finite-horizon certificates, word distortion and arithmetic scales, two-state duality, finite-bit compilation, and the global residual hierarchy. Their source bytes and loaded mathematical labels are preserved by a manifest. None is needed to repair a missing argument in the present main theorems. This separation preserves additional conclusions without assigning a referee of the present article the task of recertifying every historical derivation.

The finite-dimensional realization results are also logically independent of the repository's analytic A/B/C/D dependency graph. No raw local limit, stopped entropy recovery, global past kernel, Mosco limit, nonlinear resolvent, filtering, or posterior-contraction premise is imported as a proved step here. The standalone descriptive title reflects the mathematical results proved in this article; the repository series name serves only as a provenance identifier.
